\section{Relation to existing estimates}\label{sec:related}

Small-time heat-kernel asymptotics originate in the parametrix tradition
\cite{MinakshisundaramPleijel1949} and Varadhan's logarithmic asymptotics
\cite{Varadhan1967}.  Li--Yau estimates control positive heat-equation
solutions geometrically \cite{LiYau1986}.  These results apply far beyond the
present model, but do not normally produce the convention-specific,
finite-time array of numerical constants needed here.

The literature on logarithmic derivatives is closer to the present theorem.
Chen, Li, and Wu obtain short-time estimates for the gradient and Hessian of
the logarithmic heat kernel on complete Riemannian manifolds without curvature
restrictions \cite{ChenLiWu2023}.  Neel and Sacchelli treat arbitrary
higher-order logarithmic derivatives, uniform localized expansions, and
Morse--Bott cut loci, including standard Riemannian spheres
\cite{NeelSacchelli2025}.  Those works determine general asymptotic structure.
The result here is different in logical form: for one compact group it gives
explicit finite-$t$ majorants, computable image-tail losses, positive lower
bounds on each chart, and exact-rational numerical instances.  The
Morse--Bott interpretation of the antipode is consistent with the general
cut-locus theory.

On compact Lie groups, representation theory gives the character expansion,
and heat-kernel measures underlie the compact-group Segal--Bargmann transform
\cite{Hall1994}.  Explicit $\SU$ formulas also appear in analyses of
subelliptic heat kernels \cite{BaudoinBonnefont2008,AgrachevEtAl2009}, whose
operator and geometry differ from the Riemannian Laplacian used here.

In lattice gauge theory, the heat-kernel action of Menotti and Onofri
\cite{MenottiOnofri1981} followed Manton's geometric action \cite{Manton1980}.
Exact convolution identities are central in
two-dimensional Yang--Mills theory
\cite{GrossKingSengupta1989,Sengupta1993}.  Heat-kernel actions have also
appeared in studies of renormalization flow \cite{BitarDukeJadid1985} and in
modern work on real-time lattice actions \cite{KanwarWagman2021}.  The present
paper does not establish a flow or a continuum limit.  Its lattice
contribution is a local analytic input: explicit control of the one-plaquette
action and a precise account of what remains when plaquettes are coupled
through shared links.

The contribution is therefore the combination of explicit constants,
reproducible validation, and a complete three-chart specialization to $\SU$.
No priority claim is made for the heat-kernel
formula, the general existence of logarithmic derivative asymptotics, or the
Morse--Bott description of spherical cut loci.
